\documentclass[../DoAn.tex]{subfiles}
\begin{document}

\section{Mở đầu chương}
Chương này trình bày cơ sở lý thuyết và các phân tích liên quan đến việc xây dựng hệ thống RAG (Retrieval-Augmented Generation) Chatbot hỏi-đáp tài liệu quy định đào tạo. Nội dung chương bao gồm: phân tích bối cảnh và đặc điểm bài toán, khảo sát các nghiên cứu và hệ thống tương tự, và giới thiệu các kiến thức nền tảng như: xử lý ngôn ngữ tự nhiên (NLP), mô hình ngôn ngữ lớn (LLM), kỹ thuật Retrieval-Augmented Generation, và vector embedding.

\section{Phân tích bài toán và bối cảnh ứng dụng}

\subsection{Thực trạng tra cứu quy chế đào tạo}
Tại các trường đại học, đặc biệt là các trường có quy mô lớn như ĐHBK Hà Nội, hệ thống văn bản quy định đào tạo rất phong phú và phức tạp. Sinh viên thường gặp khó khăn trong việc:
\begin{itemize}
    \item Tìm kiếm thông tin trong các văn bản PDF dài hàng chục trang;
    \item Hiểu đúng ngữ cảnh của các điều khoản pháp lý;
    \item Xác định văn bản nào còn hiệu lực và áp dụng cho mình;
    \item Liên kết thông tin từ nhiều văn bản khác nhau.
\end{itemize}

\subsection{Mục tiêu và yêu cầu của hệ thống}
Hệ thống cần giải quyết các yêu cầu sau:
\begin{itemize}
    \item Xử lý và trích xuất nội dung từ các văn bản PDF quy chế đào tạo;
    \item Chia nhỏ văn bản thành các chunks có ngữ nghĩa, bảo toàn cấu trúc pháp lý;
    \item Lưu trữ và tìm kiếm hiệu quả bằng vector embedding;
    \item Sinh câu trả lời chính xác với trích dẫn nguồn cụ thể;
    \item Hỗ trợ tiếng Việt với độ chính xác cao.
\end{itemize}

\section{Các nghiên cứu và hệ thống liên quan}

\subsection{Sự phát triển của hệ thống hỏi-đáp}

\subsubsection{Hệ thống hỏi-đáp truyền thống}
Các hệ thống hỏi-đáp trước đây chủ yếu hoạt động theo nguyên tắc:
\begin{itemize}
    \item \textbf{Keyword matching}: Tìm kiếm dựa trên từ khóa chính xác
    \item \textbf{TF-IDF}: Đánh giá độ quan trọng của từ trong tài liệu
    \item \textbf{BM25}: Cải tiến TF-IDF cho information retrieval
\end{itemize}

\subsubsection{Hệ thống RAG hiện đại}
Ngày nay, các hệ thống RAG được phát triển dựa trên:
\begin{itemize}
    \item \textbf{Dense Retrieval}: Sử dụng neural embedding để tìm kiếm ngữ nghĩa
    \item \textbf{Large Language Models}: Mô hình ngôn ngữ lớn như GPT, Gemini, LLaMA
    \item \textbf{Retrieval-Augmented Generation}: Kết hợp retrieval và generation
    \item \textbf{Contextual understanding}: Hiểu ngữ cảnh và sinh câu trả lời mạch lạc
\end{itemize}

\subsection{Điểm mới của đề tài}
Khác với các hệ thống hiện có, đề tài này:
\begin{itemize}
    \item Xử lý đặc thù văn bản pháp lý Việt Nam với cấu trúc Chương-Điều-Khoản-Điểm;
    \item Sử dụng multi-converter pipeline để xử lý PDF có lỗi font encoding;
    \item Thiết kế chunking strategy bảo toàn metadata và ngữ cảnh pháp lý;
    \item Tích hợp trích dẫn nguồn trong câu trả lời.
\end{itemize}

\section{Kiến thức nền tảng}

\subsection{Retrieval-Augmented Generation (RAG)}

\subsubsection{Tổng quan về RAG}
RAG là kỹ thuật kết hợp retrieval (truy xuất thông tin) với generation (sinh văn bản) để tạo ra câu trả lời chính xác và có cơ sở. Pipeline RAG cơ bản gồm:
\begin{enumerate}
    \item \textbf{Indexing}: Chia nhỏ tài liệu thành chunks, tạo embedding và lưu vào vector store
    \item \textbf{Retrieval}: Tìm kiếm các chunks liên quan đến câu hỏi
    \item \textbf{Generation}: Sử dụng LLM để sinh câu trả lời dựa trên context
\end{enumerate}

\subsubsection{Ưu điểm của RAG}
\begin{itemize}
    \item Giảm hallucination của LLM bằng cách cung cấp context cụ thể
    \item Cho phép cập nhật kiến thức mà không cần retrain model
    \item Có thể trích dẫn nguồn cho câu trả lời
    \item Phù hợp với domain-specific knowledge
\end{itemize}

\subsection{Vector Embedding}

\subsubsection{Khái niệm}
Vector embedding là kỹ thuật chuyển đổi text thành vector số trong không gian nhiều chiều, sao cho các văn bản có nghĩa tương tự sẽ có vector gần nhau.

\subsubsection{E5 Multilingual Model}
Đề tài sử dụng mô hình \texttt{intfloat/multilingual-e5-large} với đặc điểm:
\begin{itemize}
    \item Dimension: 1024
    \item Hỗ trợ 100+ ngôn ngữ bao gồm tiếng Việt
    \item Hiệu suất tốt cho semantic similarity và retrieval
    \item Đã được benchmark trên nhiều tác vụ NLP
\end{itemize}

\subsection{Vector Store và FAISS}

\subsubsection{FAISS (Facebook AI Similarity Search)}
FAISS là thư viện do Meta phát triển để tìm kiếm và clustering vector hiệu quả:
\begin{itemize}
    \item Hỗ trợ nhiều loại index: Flat, IVF, HNSW
    \item Tối ưu cho GPU và CPU
    \item Có thể xử lý hàng triệu vectors
\end{itemize}

\subsubsection{IndexFlatIP}
Đề tài sử dụng \texttt{IndexFlatIP} (Inner Product) cho cosine similarity:
\begin{itemize}
    \item Exact search (không approximate)
    \item Phù hợp với dataset vừa và nhỏ (< 100K vectors)
    \item Độ chính xác cao nhất
\end{itemize}

\subsection{Large Language Models (LLM)}

\subsubsection{Gemini 2.5 Flash}
Mô hình LLM được sử dụng trong đề tài:
\begin{itemize}
    \item Phát triển bởi Google DeepMind
    \item Hỗ trợ tiếng Việt tốt
    \item Tốc độ inference nhanh
    \item Context window lớn (1M+ tokens)
    \item API dễ tích hợp
\end{itemize}

\subsection{Xử lý PDF}

\subsubsection{Thách thức với PDF tiếng Việt}
\begin{itemize}
    \item Font encoding không chuẩn gây lỗi ký tự
    \item Table extraction phức tạp
    \item Layout phức tạp với header/footer
    \item Scan PDF cần OCR
\end{itemize}

\subsubsection{Các công cụ xử lý PDF}
\begin{itemize}
    \item \textbf{Docling}: Tốt cho structure detection
    \item \textbf{PyMuPDF4LLM}: Nhanh, output markdown
    \item \textbf{olmOCR}: Vision-based, xử lý được PDF lỗi font
\end{itemize}

\section{Chunking cho văn bản pháp lý}

\subsection{Đặc điểm văn bản pháp lý Việt Nam}
Văn bản pháp lý Việt Nam có cấu trúc phân cấp rõ ràng:
\begin{verbatim}
CHƯƠNG I - QUY ĐỊNH CHUNG
    Điều 1. Phạm vi điều chỉnh
        1. Khoản 1...
        2. Khoản 2...
            a) Điểm a...
            b) Điểm b...
    Điều 2. Đối tượng áp dụng
\end{verbatim}

\subsection{Chiến lược chunking}
\begin{itemize}
    \item \textbf{Parent chunks}: 1 Điều đầy đủ + context Chương
    \item \textbf{Child chunks}: Từng Khoản trong Điều
    \item \textbf{Header chunks}: Phần đầu văn bản (căn cứ, quyết định)
    \item \textbf{Appendix chunks}: Phụ lục kèm theo
\end{itemize}

\section{Kết luận chương}
Chương 2 đã trình bày tổng quan ngữ cảnh bài toán, các nghiên cứu liên quan và nền tảng công nghệ sử dụng trong đề tài. Đặc biệt, chương đã làm rõ:
\begin{itemize}
    \item \textbf{Kiến trúc RAG}: Pipeline từ indexing đến generation
    \item \textbf{Vector embedding}: Kỹ thuật biểu diễn ngữ nghĩa văn bản
    \item \textbf{Xử lý PDF}: Các thách thức và giải pháp với tiếng Việt
    \item \textbf{Chunking pháp lý}: Chiến lược bảo toàn cấu trúc văn bản
\end{itemize}

Đây là cơ sở quan trọng để bước sang chương tiếp theo, nơi hệ thống RAG Chatbot sẽ được phân tích chi tiết và thiết kế thành các module cụ thể.

\end{document}
